% ============================================================
\chapter{Relational Algebra}
\label{ch:algebra}
% ============================================================

\section{Introduction}

When you write \texttt{SELECT name FROM students WHERE grade > 15}, you are thinking in SQL. But behind this syntax lies a mathematical language of remarkable precision: \emph{relational algebra}, invented by Edgar F.\ Codd in his landmark 1970 paper. Codd had the brilliant insight to view tables as \emph{relations} in the mathematical sense (subsets of Cartesian products) and queries as \emph{operations} on these relations: selection, projection, join, union, difference. Relational algebra is to SQL what formal logic is to natural language: a rigorous foundation that guarantees every query has a precise, unambiguous semantics.

Understanding relational algebra is understanding \emph{why} SQL works---and how the query optimiser transforms your instructions into an efficient execution plan.

\begin{definition}[Relational Algebra]
\emph{Relational algebra} is a set of operators that take one or two relations
as input and produce a relation as output. It is a \emph{procedural} language:
one specifies \emph{how} to obtain the result (sequence of operations).
\end{definition}

The operators fall into two categories:
\begin{itemize}
    \item \textbf{Set operators}: union ($\cup$), intersection ($\cap$),
          difference ($-$), Cartesian product ($\times$).
    \item \textbf{Relation-specific operators}: selection ($\sigma$), projection ($\pi$),
          join ($\bowtie$), rename ($\rho$), division ($\div$).
\end{itemize}

\section{Selection ($\sigma$)}

\begin{definition}[Selection]
The \emph{selection} $\sigma_\theta(R)$ returns the tuples of $R$ that satisfy
the condition $\theta$. The condition $\theta$ is a Boolean expression
on the attributes of $R$.
\[
\sigma_\theta(R) = \{ t \in R \mid \theta(t) = \text{true} \}
\]
\end{definition}

\begin{keyformulas}[title=\textbf{Properties of Selection}]
\begin{align}
\sigma_{\theta_1 \land \theta_2}(R) &= \sigma_{\theta_1}(\sigma_{\theta_2}(R)) \tag{cascade}\\
\sigma_{\theta_1}(\sigma_{\theta_2}(R)) &= \sigma_{\theta_2}(\sigma_{\theta_1}(R)) \tag{commutativity}\\
|\sigma_\theta(R)| &\leq |R| \tag{cardinality}
\end{align}
\end{keyformulas}

\begin{example}
Select CS students:
\[
\sigma_{\text{major} = \text{'CS'}}(\text{Students})
\]

SQL equivalent:
\begin{codeblock}[title=\textbf{SQL Equivalent of Selection}]
\begin{minted}{sql}
SELECT * FROM Students WHERE major = 'CS';
\end{minted}
\end{codeblock}

\begin{codeoutput}
\begin{verbatim}
student_id | last_name | first_name | birth_date | major
-----------+-----------+------------+------------+------
1          | Brown     | Alice      | 2004-03-15 | CS
2          | Davis     | Bob        | 2003-11-22 | CS
4          | Wilson    | David      | 2003-01-30 | CS
\end{verbatim}
\end{codeoutput}
\end{example}

\section{Projection ($\pi$)}

\begin{definition}[Projection]
The \emph{projection} $\pi_{A_1, A_2, \ldots, A_k}(R)$ returns the relation
obtained by keeping only attributes $A_1, \ldots, A_k$ of $R$ and
eliminating duplicates.
\[
\pi_{A_1, \ldots, A_k}(R) = \{ t[A_1, \ldots, A_k] \mid t \in R \}
\]
\end{definition}

\begin{example}
Project the names of students:
\[
\pi_{\text{last\_name}, \text{first\_name}}(\text{Students})
\]

\begin{codeblock}[title=\textbf{SQL Equivalent of Projection}]
\begin{minted}{sql}
SELECT DISTINCT last_name, first_name FROM Students;
\end{minted}
\end{codeblock}
\end{example}

\begin{warning}[title=\textbf{Projection and Duplicates}]
In relational algebra, projection automatically eliminates duplicates
(since a relation is a set). In SQL, you must explicitly use \texttt{DISTINCT},
because SQL works with \emph{multisets} (bags).
\end{warning}

\section{Set Operators}

\begin{definition}[Union, Intersection, Difference]
Let $R$ and $S$ be two relations with the same schema (\emph{union-compatible}):
\begin{align}
R \cup S &= \{ t \mid t \in R \lor t \in S \} \tag{union}\\
R \cap S &= \{ t \mid t \in R \land t \in S \} \tag{intersection}\\
R - S &= \{ t \mid t \in R \land t \notin S \} \tag{difference}
\end{align}
\end{definition}

\begin{theorem}[Set Properties]
\begin{itemize}
    \item $\cup$ and $\cap$ are commutative and associative.
    \item $-$ is neither commutative nor associative.
    \item $R \cap S = R - (R - S)$.
\end{itemize}
\end{theorem}

\begin{codeblock}[title=\textbf{Set Operators in SQL}]
\begin{minted}{sql}
-- Union (eliminates duplicates)
SELECT last_name FROM Students WHERE major = 'CS'
UNION
SELECT last_name FROM Students WHERE major = 'Math';

-- Intersection
SELECT student_id FROM Enrollments WHERE course_id = 101
INTERSECT
SELECT student_id FROM Enrollments WHERE course_id = 102;

-- Difference
SELECT student_id FROM Enrollments WHERE course_id = 101
EXCEPT
SELECT student_id FROM Enrollments WHERE course_id = 103;
\end{minted}
\end{codeblock}

\section{Cartesian Product ($\times$)}

\begin{definition}[Cartesian Product]
The \emph{Cartesian product} $R \times S$ produces all possible combinations
of tuples from $R$ with tuples from $S$.
\[
R \times S = \{ (t_r, t_s) \mid t_r \in R \land t_s \in S \}
\]
If $|R| = n$ and $|S| = m$, then $|R \times S| = n \times m$.
\end{definition}

\begin{warning}[title=\textbf{Cost of the Cartesian Product}]
The Cartesian product is very expensive: it multiplies the cardinalities.
In practice, it is rarely used alone; it is combined with a selection to
form a \emph{join}.
\end{warning}

\section{Join ($\bowtie$)}

\begin{definition}[Natural Join]
The \emph{natural join} $R \bowtie S$ combines tuples from $R$ and $S$ that
have equal values on their common attributes.
\[
R \bowtie S = \sigma_{\text{common attributes equal}}(R \times S)
\]
with elimination of duplicate columns.
\end{definition}

\begin{definition}[Theta-Join]
The \emph{theta-join} $R \bowtie_\theta S$ is defined as:
\[
R \bowtie_\theta S = \sigma_\theta(R \times S)
\]
where $\theta$ is any condition on the attributes of $R$ and $S$.
Special case: an \emph{equi-join} uses only equality conditions.
\end{definition}

\begin{example}
Find the names of students and the titles of courses they are enrolled in:
\[
\pi_{\text{last\_name}, \text{title}}(\text{Students} \bowtie \text{Enrollments} \bowtie \text{Courses})
\]

\begin{codeblock}[title=\textbf{Join in SQL}]
\begin{minted}{sql}
SELECT s.last_name, c.title
FROM Students s
JOIN Enrollments e ON s.student_id = e.student_id
JOIN Courses c ON e.course_id = c.course_id;
\end{minted}
\end{codeblock}

\begin{codeoutput}
\begin{verbatim}
last_name | title
----------+-----------
Brown     | Databases
Brown     | Algorithms
Davis     | Databases
Davis     | Networks
Miller    | Algorithms
Miller    | Statistics
Wilson    | Databases
Taylor    | Statistics
\end{verbatim}
\end{codeoutput}
\end{example}

\begin{proposition}[Properties of the Join]
\begin{itemize}
    \item The natural join is commutative: $R \bowtie S = S \bowtie R$.
    \item The natural join is associative: $(R \bowtie S) \bowtie T = R \bowtie (S \bowtie T)$.
    \item If $R$ and $S$ have no common attributes, $R \bowtie S = R \times S$.
\end{itemize}
\end{proposition}

\section{Rename ($\rho$)}

\begin{definition}[Rename]
The \emph{rename} operator $\rho_{B/A}(R)$ renames attribute $A$ to $B$
in relation $R$. One can also rename the relation itself: $\rho_S(R)$.
\end{definition}

Renaming is useful for:
\begin{itemize}
    \item Making two relations union-compatible.
    \item Performing a self-join (joining a relation with itself).
\end{itemize}

\begin{example}
Find pairs of students enrolled in the same course:
\[
\pi_{e_1.\text{student\_id}, e_2.\text{student\_id}}(
    \sigma_{e_1.\text{student\_id} < e_2.\text{student\_id}}(
        \rho_{e_1}(\text{Enrollments}) \bowtie_{\text{course\_id}} \rho_{e_2}(\text{Enrollments})
    )
)
\]
\end{example}

\section{Division ($\div$)}

\begin{definition}[Division]
Let $R(A, B)$ and $S(B)$ be two relations. The \emph{division} $R \div S$
returns the values of $A$ in $R$ that are associated with \emph{all}
values of $B$ in $S$.
\[
R \div S = \{ t[A] \mid t \in R \land \forall s \in S, (t[A], s[B]) \in R \}
\]
\end{definition}

\begin{theorem}[Division Expression]
Division can be expressed using other operators:
\[
R \div S = \pi_A(R) - \pi_A\big((\pi_A(R) \times S) - R\big)
\]
\end{theorem}

\begin{example}
Find students enrolled in \emph{all} courses taught by professor 2:

Let $S = \pi_{\text{course\_id}}(\sigma_{\text{prof\_id}=2}(\text{Courses}))$
(courses taught by professor 2).

$\pi_{\text{student\_id}, \text{course\_id}}(\text{Enrollments}) \div S$

gives the students enrolled in all those courses.

\begin{codeblock}[title=\textbf{Division in SQL (NOT EXISTS)}]
\begin{minted}{sql}
SELECT DISTINCT e.student_id
FROM Enrollments e
WHERE NOT EXISTS (
    SELECT c.course_id
    FROM Courses c
    WHERE c.prof_id = 2
    AND NOT EXISTS (
        SELECT 1
        FROM Enrollments e2
        WHERE e2.student_id = e.student_id
        AND e2.course_id = c.course_id
    )
);
\end{minted}
\end{codeblock}
\end{example}

\section{Query Trees and Optimization}

A relational algebra expression can be represented as a tree: leaves are
base relations, internal nodes are operators.

\begin{theorem}[Heuristic Optimization Rules]
\begin{enumerate}
    \item \textbf{Push selections down}: apply selections as early as possible
          to reduce the size of intermediate relations.
    \item \textbf{Push projections down}: project as early as possible to
          reduce the number of attributes.
    \item \textbf{Combine selection + Cartesian product into join}:
          $\sigma_\theta(R \times S) \to R \bowtie_\theta S$.
    \item \textbf{Reorder joins}: start with joins that produce the smallest
          intermediate results.
\end{enumerate}
\end{theorem}

\begin{example}
Consider the query: ``names of CS students enrolled in course 101.''

Naive expression:
\[
\pi_{\text{last\_name}}(\sigma_{\text{major='CS'} \land \text{course\_id}=101}(\text{Students} \times \text{Enrollments}))
\]

Optimized expression:
\[
\pi_{\text{last\_name}}(\sigma_{\text{major='CS'}}(\text{Students}) \bowtie \sigma_{\text{course\_id}=101}(\text{Enrollments}))
\]

The optimized version avoids the full Cartesian product and filters early.
\end{example}

\section{Algebra--SQL Correspondence Summary}

\begin{keyformulas}[title=\textbf{Relational Algebra $\leftrightarrow$ SQL Mapping}]
\begin{center}
\begin{tabular}{ll}
\toprule
\textbf{Algebra} & \textbf{SQL} \\
\midrule
$\sigma_\theta(R)$ & \texttt{SELECT * FROM R WHERE $\theta$} \\
$\pi_{A_1,\ldots,A_k}(R)$ & \texttt{SELECT DISTINCT $A_1$, \ldots, $A_k$ FROM R} \\
$R \cup S$ & \texttt{R UNION S} \\
$R \cap S$ & \texttt{R INTERSECT S} \\
$R - S$ & \texttt{R EXCEPT S} \\
$R \times S$ & \texttt{SELECT * FROM R, S} \\
$R \bowtie S$ & \texttt{SELECT * FROM R NATURAL JOIN S} \\
$R \bowtie_\theta S$ & \texttt{SELECT * FROM R JOIN S ON $\theta$} \\
$R \div S$ & Double \texttt{NOT EXISTS} \\
$\rho_{B/A}(R)$ & \texttt{SELECT A AS B FROM R} \\
\bottomrule
\end{tabular}
\end{center}
\end{keyformulas}

\section{Python Integration}

\begin{codeblock}[title=\textbf{Simulating Relational Algebra with pandas}]
\begin{minted}{python}
import pandas as pd

# Relations as DataFrames
students = pd.DataFrame({
    'student_id': [1, 2, 3, 4, 5],
    'last_name': ['Brown', 'Davis', 'Miller', 'Wilson', 'Taylor'],
    'major': ['CS', 'CS', 'Math', 'CS', 'Math']
})
enrollments = pd.DataFrame({
    'student_id': [1, 1, 2, 2, 3, 3, 4, 5],
    'course_id': [101, 102, 101, 103, 102, 104, 101, 104],
    'grade': [88.5, 76.0, 72.0, 91.0, 95.5, 82.0, 55.0, 97.0]
})

# Selection: sigma_{major='CS'}(Students)
cs = students[students['major'] == 'CS']
print("Selection:\n", cs)

# Projection: pi_{last_name}(Students)
names = students[['last_name']].drop_duplicates()
print("\nProjection:\n", names)

# Natural join
result = pd.merge(students, enrollments, on='student_id')
print("\nJoin:\n", result[['last_name', 'course_id', 'grade']])
\end{minted}
\end{codeblock}

\section*{Exercises}

\begin{exercise}
Express the following queries in relational algebra, then translate them to SQL:
\begin{enumerate}
    \item Names of professors in the ``Computer Science'' department.
    \item Titles of courses worth at least 4 credits.
    \item Names of students enrolled in the ``Databases'' course.
\end{enumerate}
\end{exercise}

\begin{exercise}
Let $R(A, B, C)$ and $S(C, D, E)$. Simplify the expression:
\[
\pi_{A, D}(\sigma_{B > 10 \land D = 'x'}(R \bowtie S))
\]
by pushing selections and projections down.
\end{exercise}

\begin{exercise}
Express the following division in SQL: ``find students enrolled in all
courses from the Computer Science department.''
\end{exercise}

\begin{exercise}
Draw the algebraic tree corresponding to the expression:
\[
\pi_{\text{last\_name}}(\sigma_{\text{grade} > 85}(\text{Students} \bowtie \text{Enrollments} \bowtie \text{Courses}))
\]
Then propose an optimized tree by applying heuristic rules.
\end{exercise}

\begin{exercise}
Prove that $R \cap S = R - (R - S)$ using the set-theoretic definitions
of the operators.
\end{exercise}
